% !TEX root = ../main.tex
\chapter{System Invariants}
\label{ch:invariants}
\label{sec:invariants}

Sections~\ref{sec:architecture} and~\ref{sec:canonical-form} argue from
architecture. This section argues from contract, which is the form a reviewer can
check without agreeing with me about mechanism.

If you were specifying a component whose assertions you intended to rely on, you
would write down invariants. Two groups are needed, and the distinction between
them is the book's practical claim in compressed form. The first group states
properties of assertion \emph{behavior}: things you would assume without thinking
about them from a database, a rule engine, or a calculation service. All of them
fail, and none can be made to hold by work on $M$. The second states properties the
\emph{system} must supply. None of them holds by default, and all of them can be
engineered.

\section{Behavioral invariants: what does not hold}

\begin{table}[ht]
  \centering
  \small
  \begin{tabular}{@{}>{\raggedright\arraybackslash}p{0.5cm}>{\raggedright\arraybackslash}p{2.9cm}>{\raggedright\arraybackslash}p{4.1cm}>{\raggedright\arraybackslash}p{1.1cm}>{\raggedright\arraybackslash}p{4.4cm}@{}}
    \toprule
    & Invariant & Statement & Holds & What breaks without it \\
    \midrule
    I1 & \textbf{Form invariance}
       & $x \sim_c x'$ implies equivalent decision-relevant outcome
       & No
       & Tests do not generalize past the wording tested \\
    \addlinespace
    I2 & \textbf{Closure}
       & $\Assert_S(p)$ and $\Assert_S(p \rightarrow q)$ imply $\Assert_S(q)$
       & No
       & Conclusions the system's own premises entail are unavailable \\
    \addlinespace
    I3 & \textbf{Decomposition}
       & $\Assert_S(p \wedge q)$ implies $\Assert_S(p)$
       & No
       & Summaries assert what their parts deny \\
    \addlinespace
    I4 & \textbf{Consistency}
       & Not both $\Assert_S(p)$ and $\Assert_S(\neg p)$
       & No
       & Two users, or one user twice, get opposite answers \\
    \addlinespace
    I5 & \textbf{Self-report fidelity}
       & Claims about what it knows track what it asserts
       & No
       & Confidence cannot be used to route or gate \\
    \addlinespace
    I6 & \textbf{Correctness}
       & $\Assert_S(p)$ implies $p$
       & No
       & Output cannot be relied on unchecked \\
    \addlinespace
    I7 & \textbf{Repeatability}
       & Same input, same session, same assertion
       & No
       & A verified answer is not a verified behavior \\
    \bottomrule
  \end{tabular}
  \caption{Behavioral invariants. I6 is the one everybody already knows fails, and
    it is the least interesting: nobody deploys these systems believing output is
    true by construction. I1 and I7 are the ones that invalidate validation.}
  \label{tab:invariants-behavior}
\end{table}

One reading to head off first. These are invariants, so \emph{holds} is a universal
claim and a single counterexample defeats it. ``No'' in the third column therefore
means the property is not guaranteed, not that behavior is always inconsistent. In
practice these systems satisfy I1--I4 on most inputs most of the time, which is
precisely why the failures are hard to catch in testing and why a specification cannot
rely on them. An invariant that holds usually is not an invariant; it is a tendency,
and nothing downstream can be built on it.

Three further observations, since the pattern matters more than the entries.

\paragraph{I6 failing is priced in; I1 failing is not.} Every deployment I have
seen assumes the system is sometimes wrong and builds review around it. Almost none
assumes that a reviewed, approved, tested behavior will not reproduce under a
rephrasing. The first is a known unknown that review handles. The second
invalidates the review.

\paragraph{I2--I4 share one cause.} These are not three independent defects. There
is no component whose function is to maintain the set of asserted claims---nothing
that closes it under consequence, decomposes conjunctions, or checks it for
contradiction. By Section~\ref{sec:what-computed} the asserted ``set'' is not a
maintained data structure at all; it is whatever successive forward passes emit.
Absent such a component, I2--I4 have nothing to be true of.

\paragraph{I5 failing has a specific and costly consequence.} A great deal of
proposed architecture routes on model-reported confidence: escalate when the model
says it is unsure, proceed when it says it is confident. I5's failure means the
routing signal is produced by the same process as the content and is not a
measurement of it. Verbalized confidence can be trained to better aggregate
calibration \citep{lin2022}, and models have some capacity to evaluate their own
outputs \citep{kadavath2022}; neither delivers the per-claim reliability a gate
requires. Aggregate calibration and per-claim boundary recognition are different
properties, and only the second supports a decision about \emph{this} answer.

That last point generalizes, and it is the shape of the whole argument: a property
that holds on average over a distribution is not a property you can gate on for an
individual case. Accuracy is such a property, as
Proposition~\ref{prop:independence} shows. So is calibration.

\section{System invariants: what must be supplied}

The constructive claim. These are conditions on $S$, not on $M$. None holds by
default; each is achievable by architecture; and together they are what must be
true for a response to cross from ordinary generative assistance into consequential
decision support.

\begin{table}[ht]
  \centering
  \small
  \begin{tabular}{@{}>{\raggedright\arraybackslash}p{0.7cm}>{\raggedright\arraybackslash}p{3.3cm}>{\raggedright\arraybackslash}p{4.6cm}>{\raggedright\arraybackslash}p{4.4cm}@{}}
    \toprule
    & Invariant & Requirement & Control that supplies it \\
    \midrule
    I8 & \textbf{Source identity}
       & Every cited source has stable identity and established authenticity
       & Controlled corpus, ingestion authentication \\
    \addlinespace
    I9 & \textbf{Temporal validity}
       & Source version and effective date are recorded and compared against $t$
       & Version metadata, corpus release pinning \\
    \addlinespace
    I10 & \textbf{Context binding}
       & $c$ is explicit before assertion, or explicitly marked unknown
       & Intake resolution, typed request objects \\
    \addlinespace
    I11 & \textbf{Applicability}
       & Evidence is filtered on scope, not similarity alone
       & Metadata-aware retrieval, scope predicates \\
    \addlinespace
    I12 & \textbf{Traceability}
       & Each consequential claim links to the evidence supporting it
       & Claim decomposition, claim-level citation \\
    \addlinespace
    I13 & \textbf{Support verification}
       & The cited evidence is checked to support the claim at the strength asserted
       & Entailment checks, policy logic, human review by tier \\
    \addlinespace
    I14 & \textbf{Reproducibility}
       & The assertion can be re-derived from recorded inputs and versions
       & Snapshot pinning, deterministic replay \\
    \addlinespace
    I15 & \textbf{Authorization}
       & Reliance is permitted at the applicable tier, with duties separated
       & Tiered policy, approval gates, access control \\
    \addlinespace
    I16 & \textbf{Reconstructability}
       & The full decision path is retained for the required period
       & Audit record, retention policy \\
    \addlinespace
    I17 & \textbf{Controlled failure}
       & Unmet conditions produce qualification, refusal, or escalation---not an
         unqualified answer
       & Output gating outside $M$ \\
    \addlinespace
    I18 & \textbf{Sufficiency}
       & The evidence set is checked for coverage of every condition the claim requires
       & Required-field or checklist verification against the rule \\
    \addlinespace
    I19 & \textbf{Defeater search}
       & A scoped search for contrary or undermining authority of equal or higher
         standing is performed and recorded
       & Governed defeater query; institution's precedence order \\
    \bottomrule
  \end{tabular}
  \caption{System invariants. I8--I9 are ordinary data governance. I10--I11 are the
    context work that Section~\ref{sec:cell-v} shows is most often absent.
    I12--I13 are what distinguishes grounding from citation. I14--I16 are what an
    audit actually asks for. I17 is what makes the others enforceable rather than
    advisory. I18--I19 correspond to the coverage and no-defeater conditions of
    Definition~\ref{def:grounding}, and exist so that every condition in the definition
    has a system obligation that establishes it.}
  \label{tab:invariants-system}
\end{table}

Two points about the second table.

\paragraph{I17 is the load-bearing one.} Without a controlled-failure path, every
other system invariant degrades into a preference. A system that detects missing
applicability and answers anyway has not enforced I11; it has measured it. And by
Section~\ref{sec:missing} the gate cannot live inside $M$, because generated
refusal is a continuation rather than a branch. I17 is therefore an architectural
requirement, not a prompt.

\paragraph{Proportionality, not maximalism.} Not every use requires all ten.
Tiering makes the controls proportional to the consequence of
reliance, and a drafting assistant needs none of I13--I19. What does not follow is
the converse: a system cannot claim auditability because it reports an accuracy
score, emits citations, or carries a generic disclaimer. Those are compatible with
failing every invariant in Table~\ref{tab:invariants-system}.

\section{Grounding Is Auditable; Accuracy Is Not}
\label{ch:auditable}
\label{sec:auditable}

\subsection{The asymmetry}

Section~\ref{sec:invariants} said what does not hold. This section says which property
can be measured instead, and why that property rather than the obvious one.

Accuracy and grounding are both properties of assertions, and they differ in a way
that decides which one a governance regime can use.

To check whether an assertion is \emph{correct}, you need to know the answer.
Sometimes you do: a deterministic
calculation can be recomputed, a figure reconciled against a system of record, a
constraint evaluated, an outcome observed later. Where correctness is cheaply decidable
it should be checked directly, and such claims
should not be routed through a language model at all.

The difficulty is the residue, and the residue is what these systems are used for:
open-world questions whose answers are not available at the moment of decision. There,
correctness can be established only against a reference built in advance---on questions
somebody already had the answer to, which is to say not the question that was asked. You cannot check correctness on the query a user just sent, because
if you could answer it you would not have sent it to the system.

To check whether an assertion is \emph{grounded}, you do not need to know the
answer. You need to know what evidence the system used and whether that evidence
is applicable at $c$. Both are properties of the system's own execution. Condition
(G1) is a check on the retrieved evidence: is it authentic, current at $t$, of
the right jurisdiction and use? Condition (G2) is a perturbation test: ablate or
contradict the evidence and see whether the assertion survives. Neither consults
a reference answer.

This has a consequence out of proportion to how obvious it sounds once stated.
Grounding can be checked \emph{at runtime, per assertion, on live traffic}.
Accuracy cannot. That is the difference between a metric and a gate. A metric
tells you, after the fact, how a sample went. A gate decides whether this
particular assertion is allowed out. Everything a regulated deployment actually
needs---blocking, escalating, qualifying, logging a reconstructable basis---requires
a gate, and only grounding supplies one.

It would be convenient to claim that grounding is stable under paraphrase where
accuracy is not, and that claim would be false. A rephrased request can retrieve
different evidence, resolve a different context, decompose into different claims, return a
different verification outcome, or engage a different precedence rule. The grounding
\emph{status} of the resulting assertion can change, and governed retrieval does not
repair form sensitivity.

What grounding supplies is different and, for a governance regime, more useful. It is
auditable \emph{per execution}. However the request was phrased, the system can inspect
the evidence, the resolved context, the verification result, the defeater search, and the
authorization path for the assertion that was actually emitted. Form-invariance failure
therefore becomes \emph{visible and attributable} rather than silent: two paraphrases that
produce different assertions produce two traces, and the traces say which condition
diverged and why. That is what accuracy cannot do, because a changed answer with no trace
is indistinguishable from a changed answer with a good reason.

So the pathology of Section~\ref{sec:canonical-form} is not repaired by measuring
grounding. It is made observable, which is the most any measurement can offer against a
property the architecture does not guarantee.

\subsection{Determinism is not the requirement}
\label{sec:determinism}

Before the measurement itself, one reading has to be blocked. The asymmetry above may
suggest that auditable systems must be deterministic, and grounding auditable because
generation is not. That reading is false and it would misdirect the remedy.

Plenty of non-deterministic computation is relied upon in regulated settings and
audited without difficulty. Monte Carlo valuation, randomized algorithms, and
sampling-based statistical procedures all produce different output on different
runs, and all are defensible. What makes them so is not determinism but three other
properties: the run is \emph{reproducible}, because the seed and the state are
recorded and the computation can be replayed; the error is \emph{bounded}, because a
convergence or confidence statement is available; and there is a
\emph{specification} saying what the computation was supposed to approximate.

Generation can satisfy the first of those, but only under conditions worth
spelling out, because they are frequently not the conditions of a commercial
deployment. Section~\ref{sec:sampling} established that the output is sampled; what
matters here is how far that can be constrained. Greedy decoding removes sampling
variance. It does not by itself produce
identical output run to run. Floating-point addition is not associative; the order in
which a reduction is performed on an accelerator depends on which kernel is selected;
and kernel selection depends on batch shape. A request batched with different
neighbours can therefore yield slightly different logits, and where two candidate
tokens are close, a different token. Sparse expert routing under capacity limits
introduces the same batch dependence more directly. Behind a hosted endpoint, the
served model version may change without notice.

Reproducibility is therefore an architectural achievement rather than a property of
the technology. It requires pinned model versions, controlled batching or
single-request inference, deterministic kernel selection, recorded decoding
parameters, and a stateless request. Those conditions are available to an
organization that self-hosts or contracts for them, and largely unavailable to one
consuming a shared public endpoint. This is why I14 appears in
Table~\ref{tab:invariants-system} as something a system must supply rather than
something it has, and a validation program that assumes replay works should test
that assumption rather than inherit it.

The point for this section is that non-determinism is not the \emph{irreducible}
gap. It can be engineered away, at a cost that is ordinary architecture.

The gap is that reproducibility over a single input buys very little when the set of
inputs you need to reason about is a semantic equivalence class you cannot
enumerate. Replaying one string tells you what that string does. It says nothing
about the materially equivalent string a different user will send tomorrow, and by
Section~\ref{sec:canonical-form} there is no bound to appeal to in its place. The
contrast with adversarial robustness is instructive: there, perturbations live in a
metric space, so a claim of the form ``behavior is stable within $\varepsilon$'' is
both meaningful and sometimes certifiable. Paraphrase admits distance measures---
embedding cosine, edit distance, entailment scores---but none with the property the
certification argument needs: that everything within $\varepsilon$ of a request is
materially equivalent to it, and everything materially equivalent is within
$\varepsilon$. Semantically identical requests can be far apart in embedding space
and materially different requests can be adjacent, which is why the corresponding
guarantee cannot be stated even approximately.

So the requirement is reproducibility plus a characterization of behavior over the
class of inputs that matter. The first is engineering and this book treats it as
such. The second is the one with no known general solution, which is why
Section~\ref{sec:auditable} argues for measuring grounding---a property of the
individual assertion and its evidence---rather than attempting to certify behavior
over a class that cannot be enumerated.

The distinguishing property is worth stating once, because the whole argument turns
on it: the systems in scope produce, as their primary output, statements that people
treat as claims about the world, and they produce them by sampling token sequences.
The first half is why the question of grounding arises at all. The second half,
developed in Section~\ref{sec:what-computed}, is why it cannot be answered by
inspecting the model.

Two boundary cases. Multimodal systems are in scope to the extent that their
text output is produced as in (1)--(4), and out of scope for their image or audio
output. Retrieval-augmented and tool-using systems built on models in scope are
squarely in scope---indeed they are the systems this book is chiefly about, since
Section~\ref{sec:measurability} argues that only such systems can be measured at
all.

Later developments---longer contexts, sparse routing, tool use, and
reinforcement training on verifiable outcomes---change capability substantially and
change none of the structural properties identified below, a claim I return to in
Section~\ref{sec:better-models}.

\subsection{What independence implies for a validation program}

Proposition~\ref{prop:independence} established that accuracy and the grounding
rate are not linked by any entailment over the feasible region. Read as a statement about
validation programs rather than about models, it says something sharper than
``accuracy is incomplete.'' It says accuracy is \emph{uninformative} about the
property that governs reliance: a reported accuracy figure is consistent with every
grounding rate, so no accuracy threshold, however high, constitutes evidence that a
system asserts within its evidence.

The asymmetry above explains why this is not symmetrical bad news. Both properties
are worth knowing, and only one of them can be established for the assertion in
front of you at the moment you need to decide about it.

\subsection{The measurement target}

The quantity to measure is the rate at which assertion outruns applicable
evidence---which, by Definition~\ref{def:hallucination}, is the hallucination rate:
\begin{equation}
  \text{hallucination rate}
  \;=_{\mathrm{def}}\;
  \Pr\bigl(\Hall(p \mid c) \;\big|\; \Assert_S(p)\bigr)
  \;=\;
  \Pr\bigl(\neg\,\Grd_S(p \mid c) \;\big|\; \Assert_S(p)\bigr).
  \label{eq:rate}
\end{equation}

This is worth pausing on, because it is where the redefinition earns its keep. Under
the standard truth-conditional definition, a hallucination rate can only be
estimated against reference answers, offline, on questions whose answers are already
known. Under Definition~\ref{def:hallucination} it is estimated from the system's own
execution: what evidence was used, and did it apply. Same word, same intuition, and a
quantity that can be computed on live traffic and used to stop an assertion. The
redefinition is what converts a retrospective score into an operable control.
This is \eqref{eq:hall} turned into a statistic. It is estimated by attribution
audit rather than comparison against reference answers, and the framework for attributing
generated statements to identified sources is established
\citep{rashkin2023}: sample assertions, and for each determine whether it is
attributable to identified evidence that is applicable at $c$.

One further decomposition is needed, because the ungrounded rate as defined pools
causes with different owners. An assertion can be ungrounded because no relevant
evidence exists in the controlled corpus; because it exists and retrieval failed to
surface it; because it was surfaced and excluded for want of applicability metadata;
because it was available and the generation path ignored it; or because policy
correctly excluded it and the system answered anyway. Those are five distinct control
failures---corpus coverage, retrieval recall, metadata completeness, gating, and policy
enforcement---and only the last two are properties of the generation path. An audit
that cannot separate them will attribute a corpus gap to the model.

Report it decomposed, not pooled. The three rows of
Table~\ref{tab:remedies} have different owners, so a single figure defeats the
purpose of having measured anything:
\begin{itemize}
  \item ungrounded rate---no evidential connection (cells ii, iv);
  \item misapplied rate---evidence present but inapplicable at $c$ (cell v);
  \item error rate---grounded and false (cell iii), reported against the corpus
        rather than the system.
\end{itemize}

\subsection{Where this can be measured, and where it cannot}
\label{sec:measurability}

One constraint is essential, and it is the point at which this book's argument
becomes an argument about architecture rather than about metrics.

The evidential base has two parts and they are not equally tractable.

For \textbf{retrieval-provided evidence}, condition (G2) is directly testable.
The evidence is identified, it is present in the context, and you can ablate it,
contradict it, or perturb it and observe whether the assertion survives. The test
is cheap, automatable, and runs on live traffic.

For \textbf{parametric evidence}---whatever the training corpus contributed to the
weights---it is not. ``Had $E$ not supported $p$'' ranges over counterfactual
training corpora, and evaluating it properly means retraining. Training-data
attribution methods approximate the counterfactual, and at frontier scale they
remain a research problem rather than an instrument. The difference between the two
cases is one of kind, not degree.

So \eqref{eq:rate} is well defined for any system and estimable only for systems
that make their evidential base explicit at inference time. An unaugmented model's
ungrounded assertion rate is a perfectly good quantity that nobody can currently
measure.

This is not a limitation of the definition. It is a requirement on the system, and
it is the strongest practical conclusion of this book. If the rate at which a
system asserts beyond its applicable evidence is the quantity that governs whether
you may rely on it, then the system has to be built so that the quantity is
observable: an identified, versioned, context-resolved evidential base at the
point of assertion, not a training corpus somewhere behind it.
Section~\ref{sec:architecture} sets out what that requires.
